% Financial Mathematics
% Exercise Template
%%%%%%%%%%%%%%%%%%%%%%%%%%%%%%%%%
\documentclass[12pt,twoside]{exam}
\usepackage{aldenexam}

\noprintanswers
%\printanswers

%%%%%%%%%%%%%%%%%%%%%%%%%%%%%%%%%%
\begin{document}

\Heading
{Financial Literacy}
{Nominal Interest Rates - Set 1}
{\textbf{Instructions: }For each problem, sketch a timeline, write an equation
  in one variable, and solve for that variable to answer the question.}

\begin{questions}

  \mathproblem[][1.5in]{
    If you put \$5 in an account earning 5\% per year compounded 5 times a year, and you leave it
    there for 5 years, how much money would you end up with?
  }

  \mathproblem[][1.5in]{
    You need \$10,000 in 3 years in order to buy a used car. If you have an account
    that earns 6\% compounded monthly, how much money do you need to deposit today to
    produce the \$10,000 by the end of 3 years?
  }

  \mathproblem[][1.5in]{
    If the population of a city grows by 3\% every quarter, how many \underline{years} (to one decimal place) will it take for
    the population to double in size?
  }

  \mathproblem[][1.5in]{
    What is the effective annual rate of interest (\% to 2 decimals) if you have a 10\% nominal annual
    rate compounded

    \bpoint Semi-annually?

    \bpoint Quarterly?

    \bpoint Monthly?

    \bpoint Daily?
  }

  \mathproblem[][1.5in]{
    If \$3,000 grows to \$4,000 over 10 years, what nominal annual rate, compounded
    monthly, did it earn?
  }

  \mathproblem[][1.5in]{
    You have an account that earns 10\% nominal annual interest, but you don’t know how
    often interest is compounded. If \$100 in this account grows to \$148.45 after 4 years, what
    was the compounding period?
  }

\end{questions}

\end{document}
